% ====================================================================
% FF3-3  합산 곡선 — 배경+4峰 전체 dQ/dV, 두 온도 (Ch1 문맥, §10 제안)
%   좌표 = eq:sum 초기 상태(평형 종 전용: L_V→0·γ=0·R_n=0)의 수치 평가
%   (gen_coords_ff3.py). U_j(T)=(−ΔH+TΔS)/F 를 온도별 재평가, w=RT/F(n=1 균일),
%   C_bg=0.05 Q_cell/V 상수 예시.
%   사용 라이브러리: arrows.meta, calc 만 (Ch1 preamble 내).
% ====================================================================
\begin{figure}[t]
\centering
\begin{tikzpicture}[x=46cm,y=0.60cm]
% ---- axes ----
\draw[-{Latex[length=1.6mm]}] (0.025,0) -- (0.025,9.6)
  node[above,font=\scriptsize] {$\dd Q/\dd V$ [$Q_\cell$/V]};
\draw[-{Latex[length=1.6mm]}] (0.025,0) -- (0.290,0)
  node[below,font=\scriptsize] {$V_n$ [V]};
\foreach \xx in {0.05,0.10,0.15,0.20,0.25}
  {\draw (\xx,0.10) -- (\xx,-0.10) node[below,font=\scriptsize] {\xx};}
\foreach \yy in {2,4,6,8}
  {\draw (0.0235,\yy) -- (0.0265,\yy) node[left,font=\scriptsize] {\yy};}
% ---- background C_bg ----
\draw[densely dotted] (0.028,0.05) -- (0.285,0.05);
\node[font=\tiny,anchor=west] at (0.255,0.42) {$C_\bg{=}0.05$};
% ---- per-transition contributions at 298.15 K (thin gray, on top of C_bg=0 baseline) ----
\draw[black!40] plot[smooth] coordinates {(0.03,1.815) (0.035,2.11) (0.04,2.432) (0.045,2.777) (0.05,3.137) (0.055,3.501) (0.06,3.854) (0.065,4.179) (0.07,4.458) (0.075,4.675) (0.08,4.814) (0.085,4.865) (0.09,4.825) (0.095,4.696) (0.1,4.488) (0.105,4.214) (0.11,3.894) (0.115,3.543) (0.12,3.18) (0.125,2.819) (0.13,2.472) (0.135,2.147) (0.14,1.849) (0.145,1.581) (0.15,1.343) (0.155,1.135) (0.16,0.9554)};
\draw[black!40] plot[smooth] coordinates {(0.05,0.5558) (0.055,0.6579) (0.06,0.7748) (0.065,0.907) (0.07,1.054) (0.075,1.215) (0.08,1.388) (0.085,1.568) (0.09,1.749) (0.095,1.926) (0.1,2.089) (0.105,2.229) (0.11,2.337) (0.115,2.407) (0.12,2.433) (0.125,2.413) (0.13,2.348) (0.135,2.244) (0.14,2.108) (0.145,1.948) (0.15,1.773) (0.155,1.591) (0.16,1.411) (0.165,1.237) (0.17,1.074) (0.175,0.9251) (0.18,0.791) (0.185,0.6721) (0.19,0.5682) (0.195,0.4782)};
\draw[black!40] plot[smooth] coordinates {(0.065,0.2276) (0.07,0.2705) (0.075,0.3201) (0.08,0.3768) (0.085,0.4408) (0.09,0.512) (0.095,0.5898) (0.1,0.673) (0.105,0.7595) (0.11,0.8467) (0.115,0.931) (0.12,1.008) (0.125,1.075) (0.13,1.125) (0.135,1.157) (0.14,1.168) (0.145,1.156) (0.15,1.124) (0.155,1.073) (0.16,1.006) (0.165,0.9283) (0.17,0.8439) (0.175,0.7566) (0.18,0.6702) (0.185,0.5872) (0.19,0.5096) (0.195,0.4386) (0.2,0.3748) (0.205,0.3183) (0.21,0.269)};
\draw[black!40] plot[smooth] coordinates {(0.14,0.2182) (0.145,0.2584) (0.15,0.3045) (0.155,0.3566) (0.16,0.4149) (0.165,0.4787) (0.17,0.5473) (0.175,0.619) (0.18,0.6918) (0.185,0.7627) (0.19,0.8286) (0.195,0.8858) (0.2,0.9307) (0.205,0.9604) (0.21,0.9728) (0.215,0.9668) (0.22,0.943) (0.225,0.9031) (0.23,0.8498) (0.235,0.7866) (0.24,0.717) (0.245,0.6445) (0.25,0.5721) (0.255,0.5022) (0.26,0.4366) (0.265,0.3763) (0.27,0.322) (0.275,0.2738) (0.28,0.2316)};
\node[black!40,font=\tiny] at (0.0715,5.05) {$2\!\to\!1$};
\node[black!40,font=\tiny] at (0.113,2.72) {$2\mathrm L\!\to\!2$};
\node[black!40,font=\tiny] at (0.145,1.47) {$3\!\to\!2\mathrm L$};
\node[black!40,font=\tiny] at (0.222,1.24) {$4\!\to\!3$};
% ---- total at 268.15 K (thick solid) ----
\draw[very thick] plot[smooth] coordinates {(0.03,1.675) (0.035,2.009) (0.04,2.398) (0.045,2.843) (0.05,3.343) (0.055,3.894) (0.06,4.485) (0.065,5.101) (0.07,5.72) (0.075,6.317) (0.08,6.865) (0.085,7.337) (0.09,7.712) (0.095,7.976) (0.1,8.122) (0.105,8.152) (0.11,8.071) (0.115,7.89) (0.12,7.622) (0.125,7.28) (0.13,6.878) (0.135,6.431) (0.14,5.955) (0.145,5.467) (0.15,4.981) (0.155,4.514) (0.16,4.078) (0.165,3.681) (0.17,3.33) (0.175,3.024) (0.18,2.76) (0.185,2.534) (0.19,2.336) (0.195,2.158) (0.2,1.993) (0.205,1.833) (0.21,1.676) (0.215,1.518) (0.22,1.362) (0.225,1.21) (0.23,1.063) (0.235,0.9254) (0.24,0.7987) (0.245,0.6844) (0.25,0.5831) (0.255,0.4948) (0.26,0.4189) (0.265,0.3544) (0.27,0.3) (0.275,0.2547) (0.28,0.217)};
% ---- total at 328.15 K (thick densely dashed) ----
\draw[very thick,densely dashed] plot[smooth] coordinates {(0.03,2.674) (0.035,3.034) (0.04,3.419) (0.045,3.823) (0.05,4.237) (0.055,4.652) (0.06,5.056) (0.065,5.437) (0.07,5.784) (0.075,6.084) (0.08,6.329) (0.085,6.513) (0.09,6.632) (0.095,6.686) (0.1,6.677) (0.105,6.609) (0.11,6.485) (0.115,6.313) (0.12,6.097) (0.125,5.843) (0.13,5.559) (0.135,5.252) (0.14,4.928) (0.145,4.596) (0.15,4.262) (0.155,3.935) (0.16,3.621) (0.165,3.325) (0.17,3.051) (0.175,2.803) (0.18,2.581) (0.185,2.386) (0.19,2.215) (0.195,2.066) (0.2,1.935) (0.205,1.818) (0.21,1.712) (0.215,1.612) (0.22,1.515) (0.225,1.419) (0.23,1.323) (0.235,1.225) (0.24,1.128) (0.245,1.03) (0.25,0.9346) (0.255,0.842) (0.26,0.7537) (0.265,0.6709) (0.27,0.5942) (0.275,0.5241) (0.28,0.4609)};
% ---- center-shift arrows (268 -> 328 K) ----
\draw[-{Latex[length=1.5mm]},thick] (0.0903,8.75) -- (0.0803,8.75);
\node[font=\tiny,anchor=south,align=center] at (0.0853,8.85)
  {$2\!\to\!1$: $-9.9$ mV\\($\Delta S_\rxn{<}0$)};
\draw[-{Latex[length=1.5mm]},thick] (0.2019,3.4) -- (0.2199,3.4);
\node[font=\tiny,anchor=south,align=center] at (0.211,3.5)
  {$4\!\to\!3$: $+18.0$ mV\\($\Delta S_\rxn{>}0$)};
\node[font=\tiny,anchor=west] at (0.228,3.0) {(중심 이동, $268\!\to\!328$ K)};
% ---- legend ----
\draw[very thick] (0.208,7.9) -- (0.220,7.9);
\node[font=\tiny,anchor=west] at (0.222,7.9) {총합, $268.15$ K};
\draw[very thick,densely dashed] (0.208,7.2) -- (0.220,7.2);
\node[font=\tiny,anchor=west] at (0.222,7.2) {총합, $328.15$ K};
\draw[black!40] (0.208,6.5) -- (0.220,6.5);
\node[font=\tiny,anchor=west] at (0.222,6.5) {전이별 기여, $298.15$ K};
\draw[densely dotted] (0.208,5.8) -- (0.220,5.8);
\node[font=\tiny,anchor=west] at (0.222,5.8) {배경 $C_\bg$};
\end{tikzpicture}
\caption{합산식~\eqref{eq:sum} 이 내는 전체 곡선 --- 좌표는 식 그대로의 수치 평가다: 초기 상태의 평형
종 전용 극한($L_V\!\to\!0$$\cdot$$\gamma_j{=}0$$\cdot$$R_n{=}0$)에서
$\dd Q/\dd V=C_\bg+\sum_jQ_j\,\xi_{\eq,j}(1-\xi_{\eq,j})/w_j$ 를 표~\ref{tab:staging} 초기값으로
평가했다($U_j(T)=(-\Delta H_{\rxn,j}+T\Delta S_{\rxn,j})/F$ 를 온도마다 재평가, $w_j{=}RT/F$
--- $n_j{=}1$ 균일, 배경은 상수 예시 $C_\bg{=}0.05\,Q_\cell$/V). 회색 가는 곡선이 $298.15$ K 의
전이별 기여, 굵은 실선$\cdot$파선이 $268.15/328.15$ K 의 총합이다. 온도가 오르면 (i) 폭
$w{=}RT/F$ 가 비례해 열려 봉우리가 낮고 넓어지고(면적 $Q_j$ 보존), (ii) 중심은
$\partial U_j/\partial T=\Delta S_{\rxn,j}/F$ 를 따라 \emph{전이마다 반대 방향}으로 움직인다 ---
$60$ K 동안 $2\!\to\!1$($\Delta S_\rxn{=}{-}16$ J/(mol\,K))은 $-9.9$ mV 로 왼쪽,
$4\!\to\!3$($+29$)은 $+18.0$ mV 로 오른쪽(화살표), $3\!\to\!2\mathrm L$($0$)은 제자리다. 다온도
피팅에서 온도마다 $U_j(T)$ 를 따로 읽어야 하는 까닭(\S\ref{sec:center})이 이 한 장이다.}
\label{fig:cand-ff3-3}
\end{figure}
